\section{Method}
\label{sec:method}

\subsection{Data and evaluation protocol}
MITOS-ATYPIA-14~\citep{roux2014mitos} provides $\times$20 frames of the same breast-cancer slides scanned
on an Aperio and a Hamamatsu scanner, and frame identifiers match across scanners. All training uses at
most 50 512\,px crop pairs taken from 24 frames of slide A03/H03. The pairs are
\emph{coordinate-corresponding}, not pixel-exact, so no training loss compares pixels across scanners. The
held-out slides A06, A08, A09, A13 and A16 are never used in training; they give 496 canonical 512\,px
crops (up to four per frame, tissue fraction $\geq 0.3$). For pixel metrics only, each Hamamatsu frame is
affinely registered into the Aperio grid with ECC; registration is never applied to training data.
A06 is a genuine colour outlier: its raw Aperio--Hamamatsu colour distance is 94.8 against 23.8--27.5 for
the other four slides, so every result is reported pooled (ALL), on the four \emph{typical} slides, and on
A06 alone.

\textbf{Metrics.} Colour fidelity is the LAB Wasserstein distance between an output $X$ and the
registered target $Y$, summed over channels,
\begin{equation}
W_{\mathrm{LAB}}(X,Y)=\sum_{c\in\{L,a,b\}}\int_{\mathbb{R}}\bigl|F_{X_c}(t)-F_{Y_c}(t)\bigr|\,dt ,
\label{eq:wlab}
\end{equation}
where $F$ is the empirical CDF of a channel. Because Eq.~\eqref{eq:wlab} ignores spatial layout,
histogram matching optimises it by construction, so we also report a windowed version averaged over
64$\times$64 tiles (wLAB). Structure is measured by grayscale SSIM~\citep{wang2004ssim} against the
registered target. Structural safety is the Relative Dice
$\mathrm{Dice}(B,G)/\mathrm{Dice}(A,G)$ of HoVer-Net~\citep{graham2019hovernet} nuclei on 130 Lizard
images~\citep{graham2021lizard} before ($A$) and after ($B$) normalisation against ground truth $G$, with a
pre-committed pass threshold of 0.95. HoVer-Net was first checked as an instrument: on the original images it
reaches $\mathrm{Dice}(A,G)=0.698$, which is the reference every normaliser is divided by. Diffusion training loss
is never used as a success signal.

\subsection{Ablation ladder (A0--A5)}
All rungs run img2img on a frozen SD\,1.5 at denoising strength 0.30, the best strength for four of five
slides. A0 is the frozen base; A1 adds a Canny ControlNet; A2 adds a rank-8 colour LoRA trained with the
denoising loss on Hamamatsu crops only, so direction comes from img2img; A3 combines ControlNet and LoRA;
A4 adds the generic LCM-LoRA (8 steps); A5 initialises from a LoRA warm-started on 3{,}000 unpaired
histopathology patches (MITOS, PanNuke~\citep{gamper2019pannuke} and TCGA).

\subsection{Source-conditioned translation}
A target-only LoRA learns the marginal $p(H)$, not the mapping $p(H\mid A)$. We therefore add a
source-conditioning branch: a ControlNet initialised from the SD\,1.5 U-Net encoder, with its conditioning
input widened to six channels, the source RGB image stacked with its Otsu-adaptive Canny edges. The cloned
encoder stays frozen (349M parameters); only the new conditioning encoder and the zero-initialised residual
output layers are trained (12.6M), jointly with a new rank-8 colour LoRA on the U-Net attention projections
(1.6M), so 14.2M parameters are trained in all. The loss is the standard noise-prediction loss on the latent
of the paired Hamamatsu crop, conditioned on its Aperio crop. Because the pairs are not pixel-exact, the source
is shifted by a random offset of up to 8\,px at each step, so the model cannot rely on exact alignment. Of the
50 pairs, 39 train the model and 11, from five whole frames, form a validation set; training uses AdamW at
learning rate $10^{-4}$, batch size 1 and bf16 for 4{,}000 steps (25\,min on one RTX~3090), and the
checkpoint with the lowest validation loss (step 3{,}750) is kept. Before this run the recipe had to overfit an
8-pair subset: it failed at 300 steps, before the zero-initialised layers had moved, and passed at 2{,}000.
The \HtoA{} model is trained the same way with the roles swapped. Shuffled- and zero-source controls confirm
that the model uses the source (SSIM 0.147 with the correct source against 0.047 with a shuffled one, in
training-domain control runs).

At inference, random img2img noise is replaced by deterministic DDIM inversion~\citep{song2021ddim}: the
source latent is inverted for 50 steps (guidance 1.0, source-conditioned) and then denoised for 50 DDIM
steps with guidance 2.0 under the same source condition (Fig.~\ref{fig:pipeline}).

\begin{figure}[t]
\centering
\begin{tikzpicture}[
  node distance=3.2mm and 2.6mm,
  box/.style={draw, rounded corners=1.5pt, align=center, font=\scriptsize, inner sep=2.2pt, minimum height=6mm},
  io/.style={box, fill=gray!12},
  ours/.style={box, fill=blue!10},
  arr/.style={-{Latex[length=1.6mm]}, semithick}]
\node[io] (a) {Aperio\\crop $A$};
\node[box, right=of a] (e) {VAE\\encoder $E$};
\node[ours, right=of e] (inv) {DDIM\\inversion};
\node[ours, right=of inv] (fwd) {DDIM\\denoising};
\node[box, right=of fwd] (d) {VAE\\decoder $D$};
\path (inv.north) -- (fwd.north) coordinate[midway] (mid);
\node[ours, above=5mm of mid] (cn) {source ControlNet + colour LoRA};
\node[ours, below=6mm of d] (f3) {F3 colour\\fusion};
\node[io, left=6mm of f3] (out) {normalised\\output};
\draw[arr] (a) -- (e);
\draw[arr] (e) -- (inv);
\draw[arr] (inv) -- (fwd);
\draw[arr] (fwd) -- (d);
\draw[arr] (cn.south -| inv.north) -- (inv.north);
\draw[arr] (cn.south -| fwd.north) -- (fwd.north);
\draw[arr, dashed] (a.north) |- (cn.west);
\draw[arr] (d) -- node[right, font=\tiny]{$H_{\mathrm{pred}}$} (f3);
\draw[arr] (f3) -- (out);
\draw[arr, dashed] (a.south) |- node[pos=0.75, below, font=\tiny]{raw $A$ and $A_V=D(E(A))$} ([yshift=-1.5mm]f3.south west) -- ([yshift=-1.5mm]f3.south west -| f3.south) -- (f3.south);
\end{tikzpicture}
\caption{Final SD\,1.5 pipeline. The source conditions the ControlNet during inversion and denoising
(dashed), and the untouched source supplies fine structure to the F3 fusion. Shaded blue: components
trained or added in this work; the SD\,1.5 U-Net and VAE encoder stay frozen.}
\label{fig:pipeline}
\end{figure}

\subsection{The VAE ceiling, fusion and decoder fine-tuning}
Every diffusion output is decoded by the VAE, which compresses a 512$\times$512$\times$3 crop 48$\times$
into a 64$\times$64$\times$4 latent. We measure the ceiling this imposes by self-reconstruction,
$\mathrm{SSIM}(x, D(E(x)))$ with a deterministic encode, and test three ways around it.

\textbf{Decoder swap.} The pretrained \texttt{sd-vae-ft-mse} decoder shares the stock encoder bit for bit,
so it decodes the same latents without retraining anything upstream.

\textbf{Detail fusion.} Instead of accepting the VAE-resynthesised image, we move only the learned
colour change back onto the untouched source. All arithmetic is in floating-point CIELAB; we write
$X=(L_X,a_X,b_X)$ for an image in LAB, $H=H_{\mathrm{pred}}$ for the translated prediction and $G_\sigma$ for
a Gaussian blur of width $\sigma$ pixels. Three variants were compared. F1 keeps the source luminance and
takes chroma from the prediction,
\begin{equation}
F_1=(L_A,\;a_H,\;b_H).
\label{eq:f1}
\end{equation}
F2 keeps the prediction's low-frequency luminance and re-injects the source's high-frequency detail,
\begin{equation}
F_2=\bigl(G_\sigma(L_H)+\alpha\,[L_A-G_\sigma(L_A)],\;a_H,\;b_H\bigr).
\label{eq:f2}
\end{equation}
F3 transfers all three channels, but only as a smoothed residual measured against the model's own
reconstruction of the source, $A_V=D(E(A))$:
\begin{equation}
F_3=A+\beta\,G_\sigma\!\bigl(H-A_V\bigr).
\label{eq:f3}
\end{equation}
Each term of Eq.~\eqref{eq:f3} has one job. Subtracting $A_V$ cancels what the VAE does to every image
(its blur and colour offset), so $H-A_V$ approximates the scanner shift alone; $G_\sigma$ removes
high-frequency texture that the generator invented; and $\beta$ sets how much of the shift is applied.
Because the added residual is smooth, F3 can change only low-frequency colour, and the fine structure of
$F_3$ is that of $A$ by construction. Variant and parameters were chosen on internal training-domain crops
and frozen ($\sigma=8$, $\beta=0.5$) before any held-out run. The real Hamamatsu image is never read to
construct a fused output; one construction bug that briefly did so (it produced SSIM $\approx0.998$) was
caught, and the corrected run was checked against the raw source.

\textbf{Decoder fine-tune.} With the encoder frozen, so that the latent space and every trained component
remain valid, we fine-tune only the decoder (49.5M parameters, initialised from \texttt{ft-mse}) on about
2{,}700 non-held-out H\&E patches (MITOS slides of both scanners, TCGA, PanNuke and the A03/H03 crops)
with an L1 + 0.3\,LPIPS~\citep{zhang2018lpips} loss for 3{,}000 steps. Checkpoints are selected on whole
non-held-out validation slides.

\subsection{Baselines}
\label{sec:baselines}
Three classical normalisers (Macenko, Reinhard, histogram matching) use one reference image each. Three
learned normalisers are trained on exactly the same 50 A03/H03 pairs as the colour LoRA, so every trained
method has the same data budget. StainNet~\citep{kang2021stainnet} is a three-layer 1$\times$1 convolutional
colour map (1{,}283 parameters) trained with a paired L1 loss; ParamNet~\citep{kang2024paramnet} uses a
ResNet18 to predict the weights of a per-image 1$\times$1 colour map, trained with its CycleGAN-style recipe;
StainGAN~\citep{shaban2019staingan} is a standard CycleGAN~\citep{zhu2017cyclegan} (ResNet-9 generator,
PatchGAN discriminator), trained here because its authors released no weights. Public MITOS-trained
checkpoints were deliberately not used, because their training splits are unknown and could include our
held-out slides. All baselines use final-step weights (no held-out checkpoint selection) and are scored by
the same crop and registration pipeline as the diffusion methods.

\subsection{SDXL transfer and exploratory probes}
The best SD\,1.5 design is ported to SDXL~\citep{podell2023sdxl} with its own rank-8 LoRA, a
source-conditioned ControlNet and an fp32 VAE (SDXL's VAE produces NaNs in fp16), at 512\,px and natively at
1024\,px, without repeating the ablation ladder. Because the 1024\,px model is data-starved at 50 pairs, two
supplementary runs use 96 non-overlapping and 144 overlapping crop pairs from the same A03/H03 frames.

Two newer backbones were probed for feasibility only. SD\,3.5-large, a transformer-based latent model with a
16-channel VAE, was tested for LoRA training cost and colour transfer. PixCell, a 0.6B-parameter
histopathology diffusion transformer that uses the same 16-channel VAE and is conditioned on a single UNI2-h
image embedding rather than text, was tested for cost, reconstruction ceiling and zero-fine-tuning
generation. PixCell is outside the proposal's scope and is reported as exploration only.
